\section{Introduction}
Pooled CRISPR screens with single-cell RNA readout (Perturb-seq) measure how the transcriptome responds to perturbing a gene \citep{dixit2016,adamson2016,norman2019,replogle2022}. Because the space of possible perturbations and combinations is far larger than any screen, many methods aim to predict responses to perturbations that have not been measured, ranging from graph neural networks \citep{roohani2024} and latent-variable models \citep{lotfollahi2023} to fine-tuned single-cell foundation models \citep{cui2024}. If accurate, such predictions could prioritise which experiments to run.

Several recent evaluations question how much of this is achieved. Simple baselines (the mean of training perturbations, an additive model for combinations, or a linear model) match or beat deep-learning methods \citep{ahlmanneltze2025,csendes2025,wu2024perturbench,li2024comparison}. Part of the explanation is \emph{systematic variation}: most perturbations shift expression in a shared direction relative to controls, so metrics computed against the control state reward predicting that shared shift \citep{vinastorne2025}. In turn, \citet{miller2026} argued that common metrics, mean squared error in particular, are poorly calibrated and favour uninformative baselines, and that deep models do outperform them under calibrated metrics. Whether unseen perturbations are predictable from available information therefore remains contested, and the answer depends on how evaluations are designed.

Three design questions recur. First, how much of a held-out response is reproducible at all: a model cannot be expected to exceed the agreement between two independent measurements of the same perturbation. Second, are training and evaluation targets statistically independent? For example, responses are usually computed against a single shared control population, whose noise is common to every response. Third, are conclusions specified before seeing the data they are tested on, or tuned on the same benchmark?

Existing evaluations answer parts of these questions. Systema \citep{vinastorne2025} replaces the control reference with the centroid of perturbed cells, so that predicting the shared shift earns no credit. \citet{csendes2025} noted that common benchmark datasets have low perturbation-specific variance. \citet{miller2026} calibrated metrics against positive and negative controls, including a technical duplicate. Our protocol combines these ideas and adds three elements: control cells split between training and evaluation targets; reproducibility limits computed from disjoint batches rather than from random cell splits; and a calibration check against more than one uninformative prediction. It is also applied to several screens under an analysis plan fixed in advance, so that conclusions from a development dataset are tested rather than assumed.

Here we address these questions directly. We built an evaluation protocol with batch-matched pseudobulk responses, disjoint control cells for training and evaluation targets, perturbation-specific metrics, batch-disjoint reproducibility limits and explicit metric-calibration controls. We developed models and baselines on one discovery dataset \citep{norman2019}. We then fixed an analysis plan, with inclusion rules, models, metrics and Holm-corrected confirmatory tests, before analysing further screens. Four of six candidates met the pre-specified inclusion rules, spanning CRISPRa and CRISPRi, three cell types and two laboratories \citep{replogle2022,tian2021}. We ask whether learned models beat the mean of other perturbations for genes absent from training, whether gene representations carry information about those genes, how performance changes with the number of training perturbations, and whether a measurement of the same perturbation in another cell line is informative. Finally, for gene pairs, we quantify how far combined responses depart from additivity after correcting for measurement error.
